\VerdierPhysicalPageAnchor{0057}%
\VerdierSemanticAnchor{ch1-1-5-7-continuation}
exhibits five vertical diagrams, respectively:

\VerdierSemanticAnchor{ch1-1-5-7-vertical}
\[
  T'(\sigma')\Delta(\sigma,\tau),\quad
  T'(\sigma)\Delta(\sigma',\tau),\quad
  \Delta(\sigma,\tau)T(\sigma'),\quad
  \Delta(\sigma',\tau)T(\sigma),\quad
  \Delta(\sigma+\sigma',\tau).
\]

It likewise exhibits four horizontal diagrams, respectively:

\VerdierSemanticAnchor{ch1-1-5-7-horizontal}
\[
  D(\sigma,\sigma')T(\tau),\quad
  D(\sigma',\sigma)T(\tau),\quad
  T'(\tau)D(\sigma,\sigma'),\quad
  T'(\tau)D(\sigma',\sigma).
\]

\VerdierSemanticAnchor{ch1-1-5-7-conclusion}
By the induction hypothesis, the diagrams $D(\sigma,\sigma')$ and
$D(\sigma',\sigma)$ are respectively $\alpha(\sigma,\sigma')$- and
$\alpha(\sigma',\sigma)$-commutative. The first part of assertion
(1.5.4.3) follows by considering the three possible cases. The second
part of assertion (1.5.4.3) is obtained by interchanging the roles of
$\sigma$ and $\tau$. In the general case, the diagram that must be
introduced is more complicated. But, using identities (1.2.1.3), the
argument reduces to that given above.

\VerdierSemanticAnchor{ch1-1-5-8}
\noindent\textbf{Remark 1.5.8.}\quad
Proposition (1.5.2) applies in particular when the cocycle $\alpha$ is
trivial. One then checks that the proposition and its proof extend
\emph{verbatim} to the case where the category $\mathcal C'$ is no
longer necessarily additive.

\VerdierSemanticAnchor{ch1-1-5-9}
\noindent\textbf{1.5.9.}\quad
Let

\VerdierSemanticAnchor{ch1-1-5-9-functors}
\[
  (F_1,p_1),(F_2,p_2)\colon
  (\mathcal C,\mathbb Z^n,T,c)
  \longrightarrow
  (\mathcal C',\mathbb Z^n,T',c')
\]
be two $\mathbb Z^n$-functors twisted respectively by two normalized
cocycles $\alpha_1,\alpha_2\in
Z^2(\mathbb Z^n,\mathbb Z/2\mathbb Z)$, and let
$\beta\colon\mathbb Z^n\to\mathbb Z/2\mathbb Z$ be a map such that

\VerdierSemanticAnchor{ch1-1-5-9-eq1}
\begin{equation}
  \alpha_2(\sigma_1,\sigma_2)\alpha_1^{-1}(\sigma_1,\sigma_2)
  =\beta(\sigma_2)\beta(\sigma_1\sigma_2)^{-1}\beta(\sigma_1),
  \qquad \sigma_1,\sigma_2\in\mathbb Z^n,
  \tag{1.5.9.1}
\end{equation}
and let $u\colon(F_1,p_1)\to(F_2,p_2)$ be a $\beta$-morphism of
functors (1.4.6). Let $e_1,\ldots,e_n$ be the canonical basis of
$\mathbb Z^n$. For every $i\in[1,n]$, the following diagram is
$\beta(e_i)$-commutative:

\VerdierSemanticAnchor{ch1-1-5-9-eq2}
\begin{equation}
  \resizebox{0.78\textwidth}{!}{%
    $\begin{tikzcd}[ampersand replacement=\&,column sep=6em,row sep=4em]
      F_1T(e_i)
        \arrow[r,"{p_1(e_i)}"]
        \arrow[d,"{u*T(e_i)}"']
        \& T'(e_i)F_1
          \arrow[d,"{T'(e_i)*u}"] \\
      F_2T(e_i)
        \arrow[r,"{p_2(e_i)}"]
        \& T'(e_i)F_2
    \end{tikzcd}$%
  }\;.
  \tag{1.5.9.2}
\end{equation}
